% from main.tex
%%%%%%%%%%  scan idx 424  |  folio 411  %%%%%%%%%%
\[
  x\leadsto f^*(x)\leadsto f^*(x)x'\leadsto f_*(f^*(x)x')
\]
est
injectif, donc $f^*$ l'est aussi.

\textbf{6.5.} Soit $X$ un schéma noethérien de dimension $d$, ayant un faisceau ample $L$. On munit $K^{\bullet}(X)$ de la filtration définie comme suit (cf. [1]): $x$ appartient à $K^{\bullet}(X)_i$ si et seulement si pour tout fermé $Y$ de $X$, il existe un complexe $M^\bullet(Y)$, parfait, tel que:
\[
  \text{i) } \cl(M^\bullet(Y))=x;
\]
\[
  \text{ii) } \operatorname{Codim}(\Supp(H^*(M^\bullet(Y)))\cap Y,Y)\ge i,
  \quad\text{avec } H^*(M^\bullet(Y))=\bigoplus_{i\ge0}H^i(M^\bullet(Y)).
\]
Les $K^{\bullet}(X)_i$, ainsi définis, sont des sous-groupes additifs de $K^{\bullet}(X)$. En effet, soient $x,x'\in K^{\bullet}(X)_i$, et pour un fermé $Y$ donné dans $X$, $M^\bullet$ et $M'{}^\bullet$ les complexes correspondants; alors $x+x'$ est la classe du complexe $M^\bullet\oplus M'{}^\bullet$, et
\[
\Supp(H^*(M^\bullet\oplus M'{}^\bullet))
 =\Supp(H^*(M^\bullet)\oplus H^*(M'{}^\bullet))
 =\Supp(H^*(M^\bullet))\cup\Supp(H^*(M'{}^\bullet));
\]
comme la codimension d'une réunion finie de fermés est la borne inférieure des codimensions des fermés, on a bien $x+x'\in K^{\bullet}(X)_i$.

\textbf{Proposition 6.6.} La filtration 5.5 de $K^{\bullet}(X)$ est une filtration d'anneau et $K^{\bullet}(X)_{d+1}=0$.

Soient $x\in K^{\bullet}(X)_i$, $x'\in K^{\bullet}(X)_j$, et $Y$ un fermé de $X$. Il existe un complexe $M^\bullet$ tel que $x=\cl(M^\bullet)$, et $\codim(\Supp(H^*(M^\bullet))\cap Y,Y)\ge i$; il existe un complexe $M'{}^\bullet$ tel que $x'=\cl(M'{}^\bullet)$, et :
\[
\codim\bigl(\Supp(H^*(M^\bullet))\cap\Supp(H^*(M'{}^\bullet))\cap Y,
            \Supp(H^*(M^\bullet))\cap Y\bigr)\ge j.
\]
On a alors : $xx'=\cl(M^\bullet\overset{\mathbf L}{\otimes}M'{}^\bullet)$, et :
\[
\codim\bigl(\Supp(H^*(M^\bullet))\cap\Supp(H^*(M'{}^\bullet))\cap Y,Y\bigr)\ge i+j
\]
(EGA $0_{\mathrm{IV}}$ 14.2.2.3). Il
